
\section{数值实验：纹理、密度平台、幂律谱与分形维数}

\subsection{实现与数据产物（可复现材料）}

我们实现了可运行的无损 HPA-CA 原型（含 uplift 记录与可逆回放）。实现与最小复现入口见附录~\ref{app:reproducibility}，核心脚本位于本论文目录的 \texttt{scripts/hpa\_ca\_lossless.py}，其输出包括：
\begin{itemize}
  \item 可见层时空图（$T\times L$ 的 $0/1$ 占据矩阵渲染）；
  \item uplift 日志图（$T\times (L/6)$ 的离散 uplift 码）；
  \item 密度曲线与密度功率谱；
  \item 盒计数曲线与估计分形维数；
  \item 原始数组数据（用于再分析）。
\end{itemize}

为便于“可缓存、可复现”的实验组织，本论文采用内容寻址的产物目录与自动生成的 \LaTeX\ 片段。运行 \texttt{scripts/run\_all.py} 将生成 \texttt{artifacts/} 下的带 \texttt{manifest.json} 的运行目录，并在 \texttt{sections/generated/} 下生成可直接 \texttt{\textbackslash input} 的表格片段，例如：
$$
\texttt{\textbackslash input\{sections/generated/hpa\_ca\_example\_run\_summary\}}
$$

\subsection{实验设置}

实验使用周期边界的一维环形链。典型设置如下：
\begin{itemize}
  \item 空间长度：$L=300$（$L$ 为 $6$ 的倍数）
  \item 时间步数：$T=200$
  \item 初态：
  \begin{itemize}
    \item 随机 Bernoulli$(p)$（扫描 $p\in\{0.1,\dots,0.9\}$，多随机种子）
    \item 单激发（仅一位为 $1$）用于可视化“波传播/纹理生成”
  \end{itemize}
\end{itemize}

度量指标包括：
\begin{itemize}
  \item 密度：$\rho_t=\frac{1}{L}\sum_{i=1}^L s_t(i)$；
  \item uplift 比例：$\mathrm{frac}(u=\delta)$（$\delta\in U$）；
  \item 密度时间序列功率谱幂律拟合斜率 $\alpha$（log-log 线性拟合）；
  \item 时空占据图的盒计数分形维数 $D$（简单 box-counting 估计）。
\end{itemize}

\subsection{观测量定义与估计方法（可审计）}

本节中关于 $\alpha$ 与 $D$ 的讨论严格视为\textbf{数值观测}，并配套可复现的估计器与误差指标（以便审计与复核），而不将其表述为已证定理。

\paragraph{功率谱幂律斜率（$\alpha$）。}
对密度序列 $\rho_t$ 去均值，取离散 Fourier 变换得到 periodogram 形式的功率谱估计 $\mathrm{PSD}(f)$。在某个频带 $f\in[f_{\min},f_{\max}]$ 内做 log-log 线性回归：
$$
\log \mathrm{PSD}(f)\approx -\alpha\log f + c,
$$
并将回归斜率的标准误作为 $\alpha$ 的不确定性指标（该误差反映拟合残差，而不等同于严格的统计置信界）。

\paragraph{盒计数维数（$D$）。}
把可见层时空图视为二值图像 $I(t,i)=s_t(i)$。对盒尺度 $\varepsilon$（取若干 $2^k$）计算占用盒数 $N(\varepsilon)$，在 log-log 坐标做线性回归：
$$
\log N(\varepsilon)\approx D\log(1/\varepsilon)+c,
$$
并同样记录回归斜率的标准误与 $R^2$。必要时可引入 burn-in（丢弃早期时间步）以隔离初始化影响。

\paragraph{误差条（跨 seed）。}
对固定 $p$ 的多随机种子样本，采用样本均值与正态近似的 $95\%$ 置信区间（$1.96\cdot \mathrm{SE}$）作为误差条；其目的在于提供可比较的粗粒度稳定性指标，而非宣称严格分布结论。

\subsection{边界情形与一致性检查}

作为最低门槛的一致性审计，我们对若干边界/极端初态进行：\textbf{(i) uplift 回放的逐步可逆性检查}，以及 \textbf{(ii) 尾段周期性检测}（若末段严格满足周期 $p$，记录 $p$；否则记为 $0$）。结果片段由脚本自动生成如下：

\begin{tabular}{lllrrrllr}%
\hline%
\textbf{case}&\textbf{kind}&\textbf{L}&\textbf{T}&\textbf{seed}&\textbf{p}&\textbf{invert}&\textbf{tail\_period}&\textbf{final\_rho}\\%
\hline%
all\_zeros&all\_zeros&300&200&1&0.5&1&1&0.000000\\%
all\_ones&all\_ones&300&200&1&0.5&1&2&0.166667\\%
single\_one&single\_one&300&200&1&0.5&1&0&0.180000\\%
single\_zero&single\_zero&300&200&1&0.5&1&0&0.206667\\%
alternating01&alternating01&300&200&1&0.5&1&1&0.000000\\%
repeat\_100001&repeat\_word6&300&200&1&0.5&1&2&0.333333\\%
bernoulli\_p&bernoulli&300&200&1&0.5&1&0&0.236667\\%
\hline%
\end{tabular}


\subsection{随机初态参数扫描（带误差条）}

对随机 Bernoulli$(p)$ 初态做 $p\times \text{seed}$ 扫描，并汇总 $\bar\rho,\bar\alpha,\bar D$ 的样本均值与 $95\%$ 置信区间（由脚本自动生成）：

\begin{tabular}{r r r l r l r l}%
\hline%
$p$&$n$&$\bar\rho$&$95\%\ \mathrm{CI}$&$\bar\alpha$&$95\%\ \mathrm{CI}$&$\bar D$&$95\%\ \mathrm{CI}$\\%
\hline%
0.10&5&0.2173&[0.1969,0.2377]&0.669&[0.291,1.047]&1.742&[1.735,1.750]\\%
0.30&5&0.1887&[0.1650,0.2124]&0.948&[0.733,1.162]&1.750&[1.748,1.753]\\%
0.50&5&0.1913&[0.1618,0.2209]&0.530&[0.412,0.648]&1.752&[1.747,1.758]\\%
0.70&5&0.1900&[0.1662,0.2138]&0.419&[0.226,0.611]&1.756&[1.749,1.762]\\%
0.90&5&0.2053&[0.1840,0.2267]&0.161&[0.077,0.246]&1.757&[1.752,1.763]\\%
\hline%
\end{tabular}


\subsection{代表性运行（自动生成片段）}

我们用一个代表性运行展示产物的组织方式与主要统计量。下表由脚本自动生成，包含运行参数、密度与盒计数维数等元数据：

\begin{tabular}{ll}%
\hline%
\textbf{key}&\textbf{value}\\%
\hline%
run\_id&\texttt{494f9f80e6fc1fc7}\\%
L&300\\%
T&200\\%
p&0.5\\%
seed&1\\%
final\_density&0.23666666666666666\\%
boxcount\_D&1.7553611550888584\\%
artifacts&\texttt{artifacts/hpa\textbackslash{}\_ca\textbackslash{}\_lossless\textbackslash{}\_run/494f9f80e6fc1fc7/}\\%
\hline%
\end{tabular}



